% =====================================================================
%  ch19_sparkroad.tex -- 第 19 章 SparkRoad-V 移植与 TD 工程操作
% =====================================================================
\section{SparkRoad-V 移植与 TD 工程操作}

本章记录把示波器从高云 GW5A 工程迁移到 \textbf{SparkRoad-V（Anlogic EG4S20BG256）}
的顶层设计、引脚约束生成方式，以及完整的 TD 工程操作步骤。

\subsection{目标板与器件}

\begin{center}
\small
\begin{tabular}{ll}
\toprule
项目 & 参数\\
\midrule
FPGA & Anlogic EG4S20BG256（Eagle 系列，BG256）\\
系统时钟 & 24\,MHz 有源晶振（K14）\\
VGA & GM7123 视频 DAC，24\,bit（RGB 各 8\,bit），HS/VS/CLK\\
ADC & AD7606 模块（8 通道 16\,bit、200\,kSPS、±5V/±10V、SPI/并口）\\
数码管 & 4 位，DIG\_1..4 + SEG\_A..H\\
按键 & 5 个独立按键（BUTTON\_L/R/U/D/C）\\
串口 & USB2UART（TX=F16，RX=E16）\\
\bottomrule
\end{tabular}
\end{center}

\subsection{顶层结构}

\file{sparkroad_scope_top.v} 的组成：

\begin{center}
\small
\begin{tabular}{lll}
\toprule
模块 & 作用 & 说明\\
\midrule
\file{uart_rx} & 接收 115200 字节 & 24\,MHz 时钟\\
\file{scope_cmd} & 示波器命令（T/R/A/L/E/S/M/C/P/H） & 产生 cfg 写脉冲\\
\file{adc_cmd} & ADC 命令（R/O/S/K/N） & 产生 CSR 写\\
\file{adc_ad7606_top} & AD7606 串行采集 + 归一化 & 10\,bit 输出\\
\file{vga_scope_top} & 示波器显示 & 640$\times$480，PIX\_DIV=1\\
\file{seg} & 数码管显示 AD 码值 & 4 位动态扫描\\
\file{uart_tx} & 每 1\,ms 上报一帧 & 16\,bit\\
\bottomrule
\end{tabular}
\end{center}

信号流：AD7606 $\rightarrow$ ADC HAL（补码$\rightarrow$偏移二进制$\rightarrow$10\,bit）
$\rightarrow$ 采集/触发/测量 $\rightarrow$ 图层合成 $\rightarrow$ VGA RGB444
$\rightarrow$ 每 4\,bit 复制为 8\,bit $\rightarrow$ GM7123。

\subsection{时钟与视频时序}

\begin{itemize}[leftmargin=2em,itemsep=2pt]
  \item 板载时钟 24\,MHz，直接作为像素时钟（\sig{PIX_DIV=1}）；
  \item 640$\times$480 行场参数：H=640+16+96+48=800，V=480+10+2+33=525；
  \item 场频 = 24\,MHz / (800$\times$525) $\approx$ 57.1\,Hz，属于显示器可接受范围；
  \item 若要严格 59.94\,Hz，可用 PLL 产生 25.175\,MHz（需要 EG4 的 PLL IP）。
\end{itemize}

\subsection{引脚与约束生成}

引脚不再手写约束，而是集中在一个 CSV：

\begin{lstlisting}[style=vsnip]
FPGA/src/adc/pins/board_pins.csv   # 信号,位宽,引脚(|分隔),说明
tools/gen_adc_pins.py              # -> FPGA_TD/eg4/SparkRoad.adc
\end{lstlisting}

CSV 覆盖 20 个信号、51 个引脚（时钟、复位、按键、串口、VGA 24\,bit、
数码管、AD7606 七根线）。修改接线只改 CSV，再运行脚本重新生成；
脚本会检查位宽匹配与引脚重复。

\subsection{TD 工程操作步骤}

\subsubsection{环境要求}

TD 必须安装 \textbf{Eagle（EG4）器件库}。若打开工程提示
\sig{Unknown device name EG4S20BG256}，说明当前 TD 只有 PH1/PH2 器件库
（本机 TD 6.0.3 即如此），需安装 TD 5.x 或向 Anlogic 索取 EG4 支持包。

\subsubsection{创建工程（Tcl）}

\begin{lstlisting}[style=vsnip]
cd <repo>/FPGA_TD/eg4
create_project -family eg4 -device EG4S20BG256 EG4S20BG256_scope .
import_device eg4_20.db -package BG256 -speed 5
add_files -list ./files.list
top_module sparkroad_scope_top
exit
\end{lstlisting}

\subsubsection{图形界面步骤}

\begin{enumerate}[leftmargin=2em,itemsep=2pt]
  \item 打开 \file{FPGA_TD/EG4S20BG256_scope.al}（或按上节重建）；
  \item 确认器件为 \sig{EG4S20BG256}，速度等级与板卡一致（默认 -5）；
  \item 源文件应为 \textbf{37 个}，顶层设为 \sig{sparkroad_scope_top}；
  \item 在约束集里加入 \file{FPGA_TD/eg4/SparkRoad.adc}；
  \item Run Synthesis $\rightarrow$ Run Place \& Route $\rightarrow$ Generate Bitstream；
  \item 用 TD Programmer 下载 \file{.bit} 到板卡。
\end{enumerate}

\subsubsection{约束文件格式}

Anlogic TD 使用 \file{.adc}，每行一条：

\begin{lstlisting}[style=vsnip]
set_pin_assignment {clk} { LOCATION = K14; IOSTANDARD = LVCMOS33; }
set_pin_assignment {vga_r[0]} { LOCATION = L5; IOSTANDARD = LVCMOS33; }
\end{lstlisting}

\subsection{上板调试顺序}

\begin{enumerate}[leftmargin=2em,itemsep=2pt]
  \item 先断开 AD7606，烧录工程，确认 VGA 出现网格/文字（证明时钟、VGA、
        图层链路正常）；
  \item 接 AD7606 的 5V/GND 与七根信号线，串口发 \sig{N1} 使能、\sig{R1} 选 ±10V；
  \item 给 IN1 输入 1\,kHz 正弦，观察波形；用 \sig{K} 调整采样间隔；
  \item 再测试触发（\sig{A1}、\sig{L}、\sig{M}）与预触发（\sig{P}）。
\end{enumerate}

\subsection{串口命令}

\begin{center}
\small
\begin{tabular}{lll}
\toprule
命令 & 作用 & 范围\\
\midrule
\sig{T<hex>} & 时基分频 & 0$\sim$FFFF\\
\sig{R<hex>} & 运行/冻结（scope） & 0/1\\
\sig{A<hex>} & 触发使能 & 0/1\\
\sig{L<hex>} & 触发电平 & 0$\sim$3FF\\
\sig{E<hex>} & 触发边沿 & 0/1\\
\sig{S<hex>} & 滚动开关 & 0/1\\
\sig{M<hex>} & 触发模式 & 0/1/2\\
\sig{C<hex>} & 触发源 & 0/1\\
\sig{P<hex>} & 预触发点数 & 0$\sim$3FF\\
\sig{H<hex>} & 保持时间 & 0$\sim$FFFF\\
\midrule
\sig{R<hex>} & ADC 量程（adc\_cmd） & 0=±5V 1=±10V\\
\sig{O<hex>} & 过采样 & 0$\sim$6\\
\sig{S<hex>} & SCLK 半周期 & 0$\sim$FFFF\\
\sig{K<hex>} & ADC 采样间隔 & 0$\sim$FFFF\\
\sig{N<hex>} & ADC 使能 & 0/1\\
\bottomrule
\end{tabular}
\end{center}

注意：\sig{R/S} 在示波器与 ADC 两个解析器中含义不同，但作用对象不同
（一个写 VGA 参数，一个写 ADC CSR），互不影响。

\subsection{已知限制}

\begin{itemize}[leftmargin=2em,itemsep=2pt]
  \item SparkRoad-V 无 DAC，CH1 暂接固定中点，仅 CH0 有效；
  \item AD7606 串行模式当前只读 DOUTA 的 ch1$\sim$ch4（ch5$\sim$8 需接 DOUTB）；
  \item 24\,MHz 像素时钟下场频约 57\,Hz，部分显示器可能需要 25.175\,MHz PLL；
  \item 本机 TD 6.0.3 缺 EG4 器件库，尚未实际综合/上板。
\end{itemize}

\subsection{验证}

\begin{center}
\small
\begin{tabular}{ll}
\toprule
项 & 结果\\
\midrule
TD 前端分析 & 37 个源文件通过\\
端口连接检查 & 35 模块 / 47 处例化 / 0 错误\\
约束生成 & 20 信号 / 51 引脚，无重复\\
综合/上板 & 待具备 EG4 器件库的 TD 环境\\
\bottomrule
\end{tabular}
\end{center}
